\documentclass[11pt,a4paper]{article}
\usepackage[utf8]{inputenc}
\usepackage[T1]{fontenc}
\usepackage[french]{babel}
\usepackage{amsmath}
\usepackage{amssymb}
\usepackage{amsthm}
\usepackage{mathtools}
\usepackage{geometry}
\usepackage{enumitem}
\usepackage{xcolor}
\usepackage{titlesec}
\usepackage{fancyhdr}
\usepackage{booktabs}
\usepackage{array}
\usepackage{longtable}
\usepackage{tcolorbox}
\usepackage{multicol}

% Configuration de la page
\geometry{margin=2.5cm}
\pagestyle{fancy}
\fancyhf{}
\fancyhead[L]{\leftmark}
\fancyhead[R]{Exercices P-value et Tests}
\fancyfoot[C]{\thepage}

% Configuration des titres
\titleformat{\section}
{\Large\bfseries\color{blue!70!black}}
{}
{0em}
{}[\titlerule]

\titleformat{\subsection}
{\large\bfseries\color{blue!50!black}}
{}
{0em}
{}

\titleformat{\subsubsection}
{\normalsize\bfseries\color{blue!30!black}}
{}
{0em}
{}

% Boîtes colorées
\newtcolorbox{astucebox}{
    colback=green!5!white,
    colframe=green!75!black,
    title=\textbf{Astuce},
    fonttitle=\bfseries
}

\newtcolorbox{rappelbox}{
    colback=blue!5!white,
    colframe=blue!75!black,
    title=\textbf{Rappel},
    fonttitle=\bfseries
}

\newtcolorbox{attentionbox}{
    colback=orange!5!white,
    colframe=orange!75!black,
    title=\textbf{Attention},
    fonttitle=\bfseries
}

\newtcolorbox{exemplebox}{
    colback=yellow!5!white,
    colframe=yellow!75!black,
    title=\textbf{Exemple},
    fonttitle=\bfseries
}

% Commandes personnalisées
\newcommand{\R}{\mathbb{R}}
\newcommand{\Var}{\text{Var}}
\newcommand{\E}{\mathbb{E}}
\newcommand{\N}{\mathcal{N}}
\newcommand{\student}{t}
\newcommand{\chideux}{\chi^2}

% Métadonnées
\title{Exercices Complets - P-value, Tests Unilatéraux/Bilatéraux\\et Lois de Probabilités}
\author{AmineGR03}
\date{\today}

\begin{document}

\maketitle

\tableofcontents
\newpage

\section{Introduction et Méthodologie}

\subsection{Objectif de ce document}

Ce document présente une série d'exercices pratiques sur :
\begin{itemize}
    \item \textbf{P-value} : interprétation et calcul
    \item \textbf{Tests unilatéraux et bilatéraux} : sur la moyenne d'un échantillon
    \item \textbf{Lois de probabilités} : normale, Student, Chi-deux
\end{itemize}

Chaque exercice comprend :
\begin{itemize}
    \item Un énoncé détaillé
    \item Des rappels théoriques nécessaires
    \item Une solution complète avec explications
    \item Des astuces et pièges à éviter
\end{itemize}

\subsection{Notations Utilisées}

\begin{itemize}
    \item $\mu$ : moyenne de la population
    \item $\sigma$ : écart-type de la population
    \item $\overline{X}$ : moyenne de l'échantillon
    \item $s$ : écart-type de l'échantillon
    \item $n$ : taille de l'échantillon
    \item $\alpha$ : niveau de signification (risque d'erreur de type I)
    \item $p$ : p-value
    \item $H_0$ : hypothèse nulle
    \item $H_1$ : hypothèse alternative
\end{itemize}

\newpage

\section{Partie 1 : P-value - Concepts et Interprétation}

\subsection{Définition et Principe}

\begin{rappelbox}
La \textbf{p-value} (valeur-p) est la probabilité d'observer une statistique de test au moins aussi extrême que celle observée, sous l'hypothèse que $H_0$ est vraie.

\begin{itemize}
    \item Si $p < \alpha$ : on rejette $H_0$ (résultat significatif)
    \item Si $p \geq \alpha$ : on ne rejette pas $H_0$ (pas de preuve contre $H_0$)
\end{itemize}
\end{rappelbox}

\subsection{Exercice 1 : Interprétation de P-value}

\textbf{Énoncé :}

Un chercheur teste l'hypothèse $H_0 : \mu = 100$ contre $H_1 : \mu \neq 100$ avec un niveau de signification $\alpha = 0.05$. Il obtient une p-value de $0.032$.

\begin{enumerate}
    \item Que peut-on conclure ?
    \item Quelle est la probabilité de commettre une erreur de type I si on rejette $H_0$ ?
    \item Si on avait choisi $\alpha = 0.01$, quelle serait la conclusion ?
\end{enumerate}

\textbf{Solution :}

\begin{enumerate}
    \item \textbf{Conclusion :}
    
    Puisque $p = 0.032 < \alpha = 0.05$, on \textbf{rejette $H_0$} au seuil de 5\%. 
    
    On conclut qu'il existe des preuves statistiques significatives que $\mu \neq 100$.
    
    \item \textbf{Probabilité d'erreur de type I :}
    
    L'erreur de type I consiste à rejeter $H_0$ alors qu'elle est vraie. Si on rejette $H_0$ avec $p = 0.032$, la probabilité de commettre cette erreur est exactement $p = 0.032 = 3.2\%$.
    
    \item \textbf{Avec $\alpha = 0.01$ :}
    
    Puisque $p = 0.032 \geq \alpha = 0.01$, on \textbf{ne rejette pas $H_0$} au seuil de 1\%.
    
    On ne peut pas conclure que $\mu \neq 100$ avec un niveau de confiance de 99\%.
\end{enumerate}

\begin{astucebox}
\textbf{Important :} La p-value n'est PAS la probabilité que $H_0$ soit vraie. C'est la probabilité d'observer des données aussi extrêmes (ou plus) sous l'hypothèse que $H_0$ est vraie.
\end{astucebox}

\subsection{Exercice 2 : Calcul de P-value (Test Bilatéral)}

\textbf{Énoncé :}

On teste $H_0 : \mu = 50$ contre $H_1 : \mu \neq 50$ avec un échantillon de taille $n = 36$, $\overline{X} = 52$, et $s = 6$. La population est supposée normale.

\begin{enumerate}
    \item Calculer la statistique de test $t$.
    \item Déterminer la p-value.
    \item Conclure au seuil $\alpha = 0.05$.
\end{enumerate}

\textbf{Solution :}

\begin{enumerate}
    \item \textbf{Statistique de test :}
    
    Comme $\sigma$ est inconnu et $n = 36$ (grand échantillon), on utilise la loi de Student avec $n-1 = 35$ degrés de liberté.
    
    $$t = \frac{\overline{X} - \mu_0}{s/\sqrt{n}} = \frac{52 - 50}{6/\sqrt{36}} = \frac{2}{6/6} = \frac{2}{1} = 2$$
    
    \item \textbf{P-value (test bilatéral) :}
    
    Pour un test bilatéral, la p-value est :
    $$p = 2 \times P(T_{35} \geq |t|) = 2 \times P(T_{35} \geq 2)$$
    
    En consultant la table de Student avec $\nu = 35$ degrés de liberté :
    \begin{itemize}
        \item $P(T_{35} \geq 2.030) \approx 0.025$
        \item $P(T_{35} \geq 1.690) \approx 0.05$
    \end{itemize}
    
    Par interpolation, $P(T_{35} \geq 2) \approx 0.027$.
    
    Donc : $$p = 2 \times 0.027 = 0.054$$
    
    \item \textbf{Conclusion :}
    
    Puisque $p = 0.054 \geq \alpha = 0.05$, on \textbf{ne rejette pas $H_0$} au seuil de 5\%.
    
    On ne peut pas conclure que $\mu \neq 50$ avec un niveau de confiance de 95\%.
\end{enumerate}

\newpage

\section{Partie 2 : Tests Unilatéraux et Bilatéraux}

\subsection{Rappels Théoriques}

\begin{rappelbox}
\textbf{Test bilatéral :}
\begin{itemize}
    \item $H_0 : \mu = \mu_0$ contre $H_1 : \mu \neq \mu_0$
    \item Zone de rejet : $|t| > t_{\alpha/2, n-1}$
    \item P-value : $p = 2 \times P(T \geq |t|)$
\end{itemize}

\textbf{Test unilatéral à droite :}
\begin{itemize}
    \item $H_0 : \mu \leq \mu_0$ contre $H_1 : \mu > \mu_0$
    \item Zone de rejet : $t > t_{\alpha, n-1}$
    \item P-value : $p = P(T \geq t)$
\end{itemize}

\textbf{Test unilatéral à gauche :}
\begin{itemize}
    \item $H_0 : \mu \geq \mu_0$ contre $H_1 : \mu < \mu_0$
    \item Zone de rejet : $t < -t_{\alpha, n-1}$
    \item P-value : $p = P(T \leq t)$
\end{itemize}
\end{rappelbox}

\subsection{Exercice 3 : Test Unilatéral à Droite}

\textbf{Énoncé :}

Un fabricant affirme que ses batteries ont une durée de vie moyenne d'au moins 120 heures. Un échantillon de 25 batteries donne $\overline{X} = 115$ heures et $s = 10$ heures. On suppose la population normale.

\begin{enumerate}
    \item Formuler les hypothèses.
    \item Calculer la statistique de test.
    \item Déterminer la p-value.
    \item Conclure au seuil $\alpha = 0.05$.
\end{enumerate}

\textbf{Solution :}

\begin{enumerate}
    \item \textbf{Hypothèses :}
    
    Le fabricant affirme "au moins 120 heures", donc :
    \begin{align*}
        H_0 &: \mu \geq 120 \\
        H_1 &: \mu < 120
    \end{align*}
    
    C'est un \textbf{test unilatéral à gauche}.
    
    \item \textbf{Statistique de test :}
    
    $$t = \frac{\overline{X} - \mu_0}{s/\sqrt{n}} = \frac{115 - 120}{10/\sqrt{25}} = \frac{-5}{10/5} = \frac{-5}{2} = -2.5$$
    
    \item \textbf{P-value :}
    
    Pour un test unilatéral à gauche :
    $$p = P(T_{24} \leq -2.5) = P(T_{24} \geq 2.5)$$
    
    En consultant la table de Student avec $\nu = 24$ :
    \begin{itemize}
        \item $P(T_{24} \geq 2.492) \approx 0.01$
        \item $P(T_{24} \geq 2.797) \approx 0.005$
    \end{itemize}
    
    Par interpolation : $p \approx 0.009$
    
    \item \textbf{Conclusion :}
    
    Puisque $p = 0.009 < \alpha = 0.05$, on \textbf{rejette $H_0$} au seuil de 5\%.
    
    On conclut que les batteries ont une durée de vie moyenne \textbf{inférieure à 120 heures}.
\end{enumerate}

\subsection{Exercice 4 : Test Bilatéral avec Variance Connue}

\textbf{Énoncé :}

On teste $H_0 : \mu = 100$ contre $H_1 : \mu \neq 100$ avec un échantillon de taille $n = 64$, $\overline{X} = 98$, et $\sigma = 8$ (connu). La population est normale.

\begin{enumerate}
    \item Calculer la statistique de test $z$.
    \item Déterminer la p-value.
    \item Conclure au seuil $\alpha = 0.05$.
    \item Donner l'intervalle de confiance à 95\% pour $\mu$.
\end{enumerate}

\textbf{Solution :}

\begin{enumerate}
    \item \textbf{Statistique de test :}
    
    Comme $\sigma$ est connu, on utilise la loi normale :
    $$z = \frac{\overline{X} - \mu_0}{\sigma/\sqrt{n}} = \frac{98 - 100}{8/\sqrt{64}} = \frac{-2}{8/8} = \frac{-2}{1} = -2$$
    
    \item \textbf{P-value (test bilatéral) :}
    
    $$p = 2 \times P(Z \geq |z|) = 2 \times P(Z \geq 2)$$
    
    En consultant la table de la loi normale :
    $$P(Z \geq 2) = 1 - \Phi(2) = 1 - 0.9772 = 0.0228$$
    
    Donc : $$p = 2 \times 0.0228 = 0.0456$$
    
    \item \textbf{Conclusion :}
    
    Puisque $p = 0.0456 < \alpha = 0.05$, on \textbf{rejette $H_0$} au seuil de 5\%.
    
    On conclut que $\mu \neq 100$ avec un niveau de confiance de 95\%.
    
    \item \textbf{Intervalle de confiance :}
    
    Pour un niveau de confiance de 95\%, $z_{\alpha/2} = z_{0.025} = 1.96$.
    
    $$IC_{95\%} = \left[\overline{X} - z_{0.025} \frac{\sigma}{\sqrt{n}}, \overline{X} + z_{0.025} \frac{\sigma}{\sqrt{n}}\right]$$
    
    $$IC_{95\%} = \left[98 - 1.96 \times \frac{8}{8}, 98 + 1.96 \times \frac{8}{8}\right] = [98 - 1.96, 98 + 1.96] = [96.04, 99.96]$$
    
    On remarque que $\mu_0 = 100$ n'est pas dans l'intervalle, ce qui confirme le rejet de $H_0$.
\end{enumerate}

\newpage

\section{Partie 3 : Lois de Probabilités}

\subsection{Loi Normale}

\begin{rappelbox}
\textbf{Loi normale $\mathcal{N}(\mu, \sigma^2)$ :}
\begin{itemize}
    \item Densité : $f(x) = \frac{1}{\sigma\sqrt{2\pi}} e^{-\frac{1}{2}\left(\frac{x-\mu}{\sigma}\right)^2}$
    \item Espérance : $\E(X) = \mu$
    \item Variance : $\Var(X) = \sigma^2$
    \item Standardisation : $Z = \frac{X - \mu}{\sigma} \sim \mathcal{N}(0, 1)$
\end{itemize}
\end{rappelbox}

\subsection{Exercice 5 : Probabilités avec la Loi Normale}

\textbf{Énoncé :}

Le poids des étudiants suit une loi normale de moyenne $\mu = 70$ kg et d'écart-type $\sigma = 8$ kg.

\begin{enumerate}
    \item Quelle est la probabilité qu'un étudiant pèse moins de 65 kg ?
    \item Quelle est la probabilité qu'un étudiant pèse entre 65 et 75 kg ?
    \item Quel est le poids tel que 90\% des étudiants pèsent moins ?
\end{enumerate}

\textbf{Solution :}

Soit $X \sim \mathcal{N}(70, 8^2)$.

\begin{enumerate}
    \item \textbf{$P(X < 65)$ :}
    
    On standardise : $Z = \frac{X - 70}{8}$
    
    $$P(X < 65) = P\left(Z < \frac{65 - 70}{8}\right) = P(Z < -0.625)$$
    
    Par symétrie : $P(Z < -0.625) = P(Z > 0.625) = 1 - \Phi(0.625)$
    
    En consultant la table : $\Phi(0.625) \approx 0.734$, donc :
    $$P(X < 65) = 1 - 0.734 = 0.266$$
    
    \item \textbf{$P(65 < X < 75)$ :}
    
    $$P(65 < X < 75) = P\left(\frac{65 - 70}{8} < Z < \frac{75 - 70}{8}\right) = P(-0.625 < Z < 0.625)$$
    
    $$= \Phi(0.625) - \Phi(-0.625) = 0.734 - 0.266 = 0.468$$
    
    \item \textbf{Quantile à 90\% :}
    
    On cherche $x$ tel que $P(X < x) = 0.90$.
    
    $$P(X < x) = P\left(Z < \frac{x - 70}{8}\right) = 0.90$$
    
    En consultant la table : $P(Z < 1.28) \approx 0.90$
    
    Donc : $\frac{x - 70}{8} = 1.28$, soit $x = 70 + 1.28 \times 8 = 70 + 10.24 = 80.24$ kg
\end{enumerate}

\subsection{Loi de Student}

\begin{rappelbox}
\textbf{Loi de Student $t_\nu$ :}
\begin{itemize}
    \item Utilisée quand $\sigma$ est inconnu et $n$ petit
    \item $\nu = n - 1$ degrés de liberté
    \item Symétrique autour de 0
    \item Quand $\nu \to \infty$, $t_\nu \to \mathcal{N}(0, 1)$
\end{itemize}
\end{rappelbox}

\subsection{Exercice 6 : Test avec Loi de Student}

\textbf{Énoncé :}

On teste $H_0 : \mu = 50$ contre $H_1 : \mu > 50$ avec un échantillon de taille $n = 16$, $\overline{X} = 52$, et $s = 4$. La population est normale.

\begin{enumerate}
    \item Calculer la statistique de test.
    \item Déterminer la valeur critique au seuil $\alpha = 0.05$.
    \item Conclure.
    \item Calculer la p-value.
\end{enumerate}

\textbf{Solution :}

\begin{enumerate}
    \item \textbf{Statistique de test :}
    
    $$t = \frac{\overline{X} - \mu_0}{s/\sqrt{n}} = \frac{52 - 50}{4/\sqrt{16}} = \frac{2}{4/4} = \frac{2}{1} = 2$$
    
    Avec $\nu = n - 1 = 15$ degrés de liberté.
    
    \item \textbf{Valeur critique :}
    
    Pour un test unilatéral à droite avec $\alpha = 0.05$ et $\nu = 15$ :
    $$t_{\alpha, \nu} = t_{0.05, 15}$$
    
    En consultant la table : $t_{0.05, 15} \approx 1.753$
    
    \item \textbf{Conclusion :}
    
    Puisque $t = 2 > t_{0.05, 15} = 1.753$, on est dans la zone de rejet.
    
    On \textbf{rejette $H_0$} au seuil de 5\%. On conclut que $\mu > 50$.
    
    \item \textbf{P-value :}
    
    $$p = P(T_{15} \geq 2)$$
    
    En consultant la table :
    \begin{itemize}
        \item $P(T_{15} \geq 1.753) = 0.05$
        \item $P(T_{15} \geq 2.131) = 0.025$
    \end{itemize}
    
    Par interpolation : $p \approx 0.032$
    
    Puisque $p = 0.032 < \alpha = 0.05$, on confirme le rejet de $H_0$.
\end{enumerate}

\subsection{Loi Chi-deux}

\begin{rappelbox}
\textbf{Loi Chi-deux $\chi^2_\nu$ :}
\begin{itemize}
    \item Utilisée pour tester la variance
    \item $\nu$ degrés de liberté
    \item Non symétrique, toujours positive
    \item Si $X_1, \ldots, X_n \sim \mathcal{N}(\mu, \sigma^2)$ i.i.d., alors :
    $$\frac{(n-1)s^2}{\sigma^2} \sim \chi^2_{n-1}$$
\end{itemize}
\end{rappelbox}

\subsection{Exercice 7 : Test sur la Variance}

\textbf{Énoncé :}

On teste $H_0 : \sigma^2 = 25$ contre $H_1 : \sigma^2 > 25$ avec un échantillon de taille $n = 20$, $s^2 = 36$. La population est normale.

\begin{enumerate}
    \item Calculer la statistique de test.
    \item Déterminer la valeur critique au seuil $\alpha = 0.05$.
    \item Conclure.
    \item Calculer la p-value.
\end{enumerate}

\textbf{Solution :}

\begin{enumerate}
    \item \textbf{Statistique de test :}
    
    $$\chi^2 = \frac{(n-1)s^2}{\sigma_0^2} = \frac{(20-1) \times 36}{25} = \frac{19 \times 36}{25} = \frac{684}{25} = 27.36$$
    
    Avec $\nu = n - 1 = 19$ degrés de liberté.
    
    \item \textbf{Valeur critique :}
    
    Pour un test unilatéral à droite avec $\alpha = 0.05$ et $\nu = 19$ :
    $$\chi^2_{\alpha, \nu} = \chi^2_{0.05, 19}$$
    
    En consultant la table : $\chi^2_{0.05, 19} \approx 30.144$
    
    \item \textbf{Conclusion :}
    
    Puisque $\chi^2 = 27.36 < \chi^2_{0.05, 19} = 30.144$, on n'est pas dans la zone de rejet.
    
    On \textbf{ne rejette pas $H_0$} au seuil de 5\%. On ne peut pas conclure que $\sigma^2 > 25$.
    
    \item \textbf{P-value :}
    
    $$p = P(\chi^2_{19} \geq 27.36)$$
    
    En consultant la table :
    \begin{itemize}
        \item $P(\chi^2_{19} \geq 27.204) = 0.10$
        \item $P(\chi^2_{19} \geq 30.144) = 0.05$
    \end{itemize}
    
    Par interpolation : $p \approx 0.09$
    
    Puisque $p = 0.09 \geq \alpha = 0.05$, on confirme qu'on ne rejette pas $H_0$.
\end{enumerate}

\newpage

\section{Partie 4 : Exercices Synthèses}

\subsection{Exercice 8 : Problème Complet}

\textbf{Énoncé :}

Un fabricant de composants électroniques affirme que la résistance moyenne de ses composants est de 100 $\Omega$ avec un écart-type de 5 $\Omega$. Un client soupçonne que la résistance moyenne est différente. Il teste un échantillon de 36 composants et obtient une résistance moyenne de 98 $\Omega$ avec un écart-type de 6 $\Omega$.

\begin{enumerate}
    \item Formuler les hypothèses appropriées.
    \item Calculer la statistique de test (on suppose $\sigma$ inconnu).
    \item Déterminer la p-value.
    \item Conclure au seuil $\alpha = 0.05$.
    \item Donner un intervalle de confiance à 95\% pour la résistance moyenne.
    \item Tester si l'écart-type est différent de 5 $\Omega$ au seuil $\alpha = 0.05$.
\end{enumerate}

\textbf{Solution :}

\begin{enumerate}
    \item \textbf{Hypothèses :}
    
    Le client soupçonne que la moyenne est différente (pas supérieure ou inférieure spécifiquement), donc :
    \begin{align*}
        H_0 &: \mu = 100 \\
        H_1 &: \mu \neq 100
    \end{align*}
    
    C'est un \textbf{test bilatéral}.
    
    \item \textbf{Statistique de test :}
    
    Comme $\sigma$ est inconnu et $n = 36$ (grand échantillon), on utilise la loi de Student :
    $$t = \frac{\overline{X} - \mu_0}{s/\sqrt{n}} = \frac{98 - 100}{6/\sqrt{36}} = \frac{-2}{6/6} = \frac{-2}{1} = -2$$
    
    Avec $\nu = n - 1 = 35$ degrés de liberté.
    
    \item \textbf{P-value :}
    
    Pour un test bilatéral :
    $$p = 2 \times P(T_{35} \geq |t|) = 2 \times P(T_{35} \geq 2)$$
    
    En consultant la table : $P(T_{35} \geq 2) \approx 0.027$
    
    Donc : $p = 2 \times 0.027 = 0.054$
    
    \item \textbf{Conclusion :}
    
    Puisque $p = 0.054 \geq \alpha = 0.05$, on \textbf{ne rejette pas $H_0$} au seuil de 5\%.
    
    On ne peut pas conclure que la résistance moyenne est différente de 100 $\Omega$ avec un niveau de confiance de 95\%.
    
    \item \textbf{Intervalle de confiance :}
    
    Pour $\nu = 35$ et $\alpha = 0.05$, $t_{0.025, 35} \approx 2.030$
    
    $$IC_{95\%} = \left[\overline{X} - t_{0.025, 35} \frac{s}{\sqrt{n}}, \overline{X} + t_{0.025, 35} \frac{s}{\sqrt{n}}\right]$$
    
    $$IC_{95\%} = \left[98 - 2.030 \times \frac{6}{6}, 98 + 2.030 \times \frac{6}{6}\right] = [98 - 2.030, 98 + 2.030] = [95.97, 100.03]$$
    
    On remarque que $\mu_0 = 100$ est dans l'intervalle, ce qui confirme qu'on ne rejette pas $H_0$.
    
    \item \textbf{Test sur l'écart-type :}
    
    Hypothèses : $H_0 : \sigma = 5$ contre $H_1 : \sigma \neq 5$ (test bilatéral)
    
    Statistique de test :
    $$\chi^2 = \frac{(n-1)s^2}{\sigma_0^2} = \frac{(36-1) \times 6^2}{5^2} = \frac{35 \times 36}{25} = \frac{1260}{25} = 50.4$$
    
    Avec $\nu = 35$ degrés de liberté.
    
    Valeurs critiques (test bilatéral) :
    \begin{itemize}
        \item $\chi^2_{0.025, 35} \approx 20.569$ (borne inférieure)
        \item $\chi^2_{0.975, 35} \approx 53.203$ (borne supérieure)
    \end{itemize}
    
    Puisque $20.569 < 50.4 < 53.203$, on n'est pas dans la zone de rejet.
    
    On \textbf{ne rejette pas $H_0$} au seuil de 5\%. On ne peut pas conclure que $\sigma \neq 5$ $\Omega$.
\end{enumerate}

\newpage

\section{Résumé et Tableaux Récapitulatifs}

\subsection{Tableau Comparatif : Tests Unilatéraux vs Bilatéraux}

\begin{table}[h]
\centering
\small
\begin{tabular}{|l|c|c|}
\hline
\textbf{Caractéristique} & \textbf{Unilatéral} & \textbf{Bilatéral} \\
\hline
Hypothèses & $H_1 : \mu > \mu_0$ ou $\mu < \mu_0$ & $H_1 : \mu \neq \mu_0$ \\
\hline
Zone de rejet & Un seul côté & Deux côtés \\
\hline
Valeur critique & $t_{\alpha, \nu}$ & $t_{\alpha/2, \nu}$ \\
\hline
P-value & $P(T \geq t)$ ou $P(T \leq t)$ & $2 \times P(T \geq |t|)$ \\
\hline
Puissance & Plus élevée & Plus faible \\
\hline
Utilisation & Direction connue & Direction inconnue \\
\hline
\end{tabular}
\caption{Comparaison des tests unilatéraux et bilatéraux}
\end{table}

\subsection{Tableau : Choix de la Loi de Probabilité}

\begin{table}[h]
\centering
\small
\begin{tabular}{|l|c|c|}
\hline
\textbf{Condition} & \textbf{Loi utilisée} & \textbf{Statistique} \\
\hline
$\sigma$ connu, $n$ quelconque & Normale $\mathcal{N}(0,1)$ & $z = \frac{\overline{X} - \mu_0}{\sigma/\sqrt{n}}$ \\
\hline
$\sigma$ inconnu, $n$ grand ($n \geq 30$) & Normale (approximation) & $z = \frac{\overline{X} - \mu_0}{s/\sqrt{n}}$ \\
\hline
$\sigma$ inconnu, $n$ petit ($n < 30$) & Student $t_{n-1}$ & $t = \frac{\overline{X} - \mu_0}{s/\sqrt{n}}$ \\
\hline
Test sur variance & Chi-deux $\chi^2_{n-1}$ & $\chi^2 = \frac{(n-1)s^2}{\sigma_0^2}$ \\
\hline
\end{tabular}
\caption{Choix de la loi selon les conditions}
\end{table}

\subsection{Formules Clés}

\begin{rappelbox}
\textbf{Statistiques de test :}
\begin{align}
z &= \frac{\overline{X} - \mu_0}{\sigma/\sqrt{n}} \quad \text{(variance connue)} \\
t &= \frac{\overline{X} - \mu_0}{s/\sqrt{n}} \quad \text{(variance inconnue)} \\
\chi^2 &= \frac{(n-1)s^2}{\sigma_0^2} \quad \text{(test sur variance)}
\end{align}

\textbf{Intervalles de confiance :}
\begin{align}
IC_{1-\alpha} &= \left[\overline{X} - z_{\alpha/2} \frac{\sigma}{\sqrt{n}}, \overline{X} + z_{\alpha/2} \frac{\sigma}{\sqrt{n}}\right] \quad \text{($\sigma$ connu)} \\
IC_{1-\alpha} &= \left[\overline{X} - t_{\alpha/2, n-1} \frac{s}{\sqrt{n}}, \overline{X} + t_{\alpha/2, n-1} \frac{s}{\sqrt{n}}\right] \quad \text{($\sigma$ inconnu)}
\end{align}
\end{rappelbox}

\subsection{Checklist pour Résoudre un Exercice}

\begin{enumerate}
    \item \textbf{Identifier le type de test :}
    \begin{itemize}
        \item Test sur la moyenne ? Variance ? Proportion ?
        \item Unilatéral ou bilatéral ?
    \end{itemize}
    
    \item \textbf{Formuler les hypothèses :}
    \begin{itemize}
        \item $H_0$ : hypothèse nulle (à tester)
        \item $H_1$ : hypothèse alternative
    \end{itemize}
    
    \item \textbf{Choisir la loi appropriée :}
    \begin{itemize}
        \item Variance connue ? $\to$ Loi normale
        \item Variance inconnue, $n$ grand ? $\to$ Loi normale (approximation)
        \item Variance inconnue, $n$ petit ? $\to$ Loi de Student
    \end{itemize}
    
    \item \textbf{Calculer la statistique de test}
    
    \item \textbf{Déterminer la p-value ou comparer à la valeur critique}
    
    \item \textbf{Conclure :}
    \begin{itemize}
        \item Si $p < \alpha$ : rejeter $H_0$
        \item Si $p \geq \alpha$ : ne pas rejeter $H_0$
    \end{itemize}
    
    \item \textbf{Interpréter dans le contexte}
\end{enumerate}

\vspace{2cm}

\begin{center}
\textbf{Fin du document}
\end{center}

\end{document}

